problem:  
Let \( M^n \subset \mathbb{H}^{n+1} \) be a complete, properly embedded minimal hypersurface whose second fundamental form satisfies     
\[
\int_M |A|^n < \infty,
\]
so that \(M\) has finite total curvature.  
Each end of \(M\) determines a subset of the ideal boundary \(S^n_\infty = \partial_\infty \mathbb{H}^{n+1}\).  

**Conjecture:**  
Every such end of \(M\) is asymptotic (in the normal exponential sense) to a minimal cone \(C_\Lambda\) in \(\mathbb{H}^{n+1}\), whose link \(\Lambda \subset S^n_\infty\) is a smooth, closed, minimal hypersurface of \(S^n_\infty\).  
Equivalently, finite total curvature forces the asymptotic boundary \(\partial_\infty M\) to be a finite union of smooth minimal hypersurfaces in \(S^n_\infty\), each corresponding to an end of \(M\).  

Why is it a "good" problem:  
This conjecture isolates a sharp and fundamental asymptotic property of minimal hypersurfaces in hyperbolic space. Existing results show partial analogues for minimal surfaces in \(\mathbb{R}^3\) (Osserman, Jorge–Meeks) and for certain classes in hyperbolic \(3\)-space, but no comprehensive theorem is known in higher dimensions.  

The problem’s depth lies in connecting three major geometric themes:  
- the *analytic decay* of the minimal surface equation in negatively curved ambient geometry,  
- the *geometric measure theory* of tangent cones and varifold limits at infinity, and  
- the *global conformal geometry* of the sphere at infinity as a boundary compactification.  

A resolution would clarify whether curvature finiteness alone controls asymptotic smoothness and rigidity in hyperbolic spaces, potentially yielding a canonical compactification theory for complete minimal hypersurfaces.  
Its formulation is concrete, precise, and expresses a natural conjectural bridge between local analytic decay and global geometric structure—qualities that make it a genuinely "good" and profound research problem.